\documentclass[11pt]{article}
\usepackage[margin=1in]{geometry}
\usepackage[hidelinks]{hyperref}
\setlength{\parindent}{0pt}
\setlength{\parskip}{0.55em}
\title{Communication and Observation Regimes for Swarm Coordination}
\author{}
\date{Working note, 24 November 2026}
\begin{document}
\maketitle

This note collects the project discussion on radio service, other modalities, and vision-only coordination. It is a working reference for selecting assumptions and comparison cases, not a comprehensive inventory, a new contract obligation, or a claim that any listed fallback guarantees safety.

\section{What reaches the decision}
Describe a regime relative to a decision: which actor can obtain which evidence, from whom, by what deadline, and with what guarantee. Record source and subject identity, meaning, observation age, error, and any acknowledged receipt. Distinguish sending, delivery, acceptance, and execution. Record clock and coordinate-frame bounds when they affect the interpretation.

The following cases are not a single ranking. A fast one-way warning and a slower acknowledged reservation may each support a decision that the other cannot.
\begin{description}
\item[Full for the maneuver.] State, intent, commitments, and receipts arrive within their required bounds. A backup path is equivalent for this purpose only if it preserves the end-to-end guarantees, including ordering or duplicate suppression where needed.
\item[Slower or older.] The message types survive but delay, jitter, and observation age grow. Bound age at use; slower motion, an earlier decision, or more physical clearance can change which evidence suffices.
\item[Lower capacity or critical traffic only.] Coarse motion envelopes or short permissions may survive without rich updates. A one-way abort with no assured receipt does not cancel a live permission.
\item[Bounded blackout.] A guaranteed delivery opportunity after a known interval can support a guarded continuation that lasts through the blackout and its detection delay.
\item[Intermittent or eventually timely.] A later contact, expected delivery, or probabilistic service need not deliver the evidence by the current deadline. Check decisions during the gap independently of renegotiation after it.
\item[One-way or selective.] Telemetry without command, command without acknowledgment, and uneven broadcast create different local histories. The usable service may differ by direction and message type.
\item[Relayed or partitioned.] Groups may coordinate internally. A time-varying relay path can exist yet arrive too late. Local replanning does not itself release another group's unresolved claim.
\item[Beacon or rendezvous only.] A presence pulse or brief contact may allow a narrow exchange. Evidence that a transmitter is alive does not establish a clear passage.
\item[No messages, live observation.] Motion, occupancy, a gate, or a shared cue can remain observable. A negative observation does not exclude an occluded approach.
\item[No live exchange, prior agreement.] Roles, reservations, maps, or schedules can still govern actions within their assumptions. Loss of the clock bound can invalidate a timetable without invalidating spatial separation.
\item[No usable shared evidence by the deadline.] Only locally justified moves and applicable prior safeguards are available. A sound response may defer entry, separate routes, or give up progress. If every safe fallback is already unavailable, the operation committed too far before communication failed.
\end{description}

Keep integrity, freshness, authorship, failure detection, and common interpretation separate from delivery. A prompt authenticated report can be stale or refer to a different map. Crews can detect the same failure at different times. On rejoin, restored transport does not settle unresolved duties or reveal what happened during the partition. Distinguish bounded silence, a still-reachable physical rendezvous that is late for the decision, and complete separation for the mission.


For fault selection, separate independent loss, burst loss, a shared relay outage, and a moving partition. Also separate corruption, stale, duplicate or reordered messages, and contradictory claims about the same duty. An authenticated speaker can still be wrong, and a physical departure from a constraint is a different fault from a false report. A clean shutdown, a bounded notification, and independently expiring timeouts give crews different evidence about which messages or permissions may still be in flight. Clock drift, mismatched plan versions, and disagreement about membership can make a delivered message unusable without any network outage.

For more background, see \href{https://dspace.mit.edu/entities/publication/108bf327-aff6-414b-a48f-b357a821c48e}{partial synchrony}, \href{https://arxiv.org/abs/1012.0009}{time-varying graphs}, and \href{https://arxiv.org/abs/cs/0006009}{knowledge in distributed systems}.

\section{Service in other modalities}
\begin{description}
\item[Optical data and discrete signals.] Two-way packets can degrade to directional data, a short decoded signal, visible light without reliable decoding, or occlusion. Distinguish a coded packet, a conventional brake light, and observed deceleration. They may come through the same camera and warrant different conclusions.
\item[Acoustic.] Data and ranging can degrade to short delayed packets, occasional pings, sound without attributable source, or no usable signal. A locatable ping need not convey intent.
\item[Wired, docked, or carried data.] Continuous exchange can degrade to a scheduled dock, delayed one-way delivery, a cached snapshot of uncertain age, or no exchange. Time to contact can matter more than link capacity.
\item[Persistent marks.] Attributed, dated records can degrade to a readable but unattributed marker, a trace of unknown age, or nothing readable. Failure to see a mark does not establish that none was left. Record what place and interval it refers to.
\item[Mechanical feedback or interlock.] A shared load can provide continuous or intermittent force or motion feedback. An interlock instead removes a conflicting action if its mechanical state is assured. These mechanisms need different failure models.
\item[Shared references and prior agreement.] Bounds on frame and clock error support a timed common plan. Landmarks, static roles, or spatial rules may still be usable when timing is not. None can convey a newly discovered fact by itself.
\end{description}

\section{Vision and local coordination}
Reliable relative tracks can support spacing and heading changes. Restricted fields of view limit this to visible neighbors. Fresh bearings or optic flow without reliable range can support some steering responses without establishing a metric stopping gap. Intermittent sightings require accounting for motion during gaps, identity switches, and observation age. A bearing may suffice for alignment, while a task tied to a particular neighbor's duty needs persistent identity. A landmark or occupancy cue may suffice for a local rule, but a patch seen clear now need not stay clear through a crossing. With no usable visual evidence, any fallback that depends on seeing a neighbor has lost its premise.

\href{https://pmc.ncbi.nlm.nih.gov/articles/PMC2234121/}{Ballerini et al.} and \href{https://www.nature.com/articles/nphys3035}{Attanasi et al.} motivate studying coordination based on local observations; the cited field work does not establish a vision-only mechanism or prescribe a controller's neighborhood size. For robot models, compare \href{https://arxiv.org/abs/2012.01245}{detection, tracking, and relative-motion estimation} with \href{https://www.nature.com/articles/s44182-025-00027-2}{direct visual projection cues}. The cited detector-and-tracker demonstration used three quadcopters, constrained motion to a horizontal plane, and supplied common migration waypoints. It does not establish a dense 3D performance bound. Detector throughput alone does not measure either safety or coordination through occlusion. Compare age of usable observations, minimum separation, and completion under the same dynamics and visibility.

\section{Applying the note}
Classify the transition as well as the destination. For example, identical blackouts before and after a rover passes its last safe stop can leave different commitments to clear the next bay. If another rover has the same local history in both cases, its policy must work in both. A visible empty bay is not a release of the first rover's reservation. A verified exit or a physical interlock that excludes conflicting entry can sometimes justify movement without reproducing the old messages. Conversely, an attributed reservation and a local occupancy check can work together if they identify the same route and time window.

For enclosed layouts, keep the travel graph, possible communication paths, sight lines, and locally known map separate. A closed passage can affect more than one of them. To isolate the information problem, start with a known layout and compare identical outages at different commitment points. Then vary the surviving modality at a fixed commitment point. Only then add an incomplete map or a shared cause of travel and communication loss. Caves are one setting for these tests alongside underground corridors and building interiors.

For a checked example, see \href{https://github.com/eternal-november/eos-research-programmes/tree/main/projects/swarm-coordination/research/communication/multimodal-bay}{SC-MULTIMODAL-HANDOFF}. The illustrative regimes above are broader than the verified model.
\end{document}
